\documentclass[10pt,a4paper]{amsart}
\usepackage[margin=19mm]{geometry}
\usepackage{amsmath,amssymb,mathtools}
\usepackage[scr=rsfs,cal=euler]{mathalfa}
\usepackage{xfrac}
\usepackage[unicode,hidelinks,bookmarks=false]{hyperref}
\newcommand{\ie}{i.e.\ }
\newcommand{\cf}{cf.\ }
\newcommand{\resp}{respectively\ }
\newcommand{\Subsection}{\subsection}
\providecommand{\nobreakdash}{\nobreak\hbox{-}}
\setlength{\emergencystretch}{2em}
\multlinegap0pt
\usepackage{mathtools}
\usepackage{amssymb}
\usepackage[normalem]{ulem}
\usepackage{cancel}
\newtheorem{proposition}{Proposition}
\newcommand{\R}{\mathbb{R}}
\title{On a Problem Posed by Brezis and Mironescu}
\author{Fanghua Lin, Malkeil Shoshan, Changyou Wang}
\date{Source: arXiv:2603.07877v1 (CC BY 4.0)}
\begin{document}
\maketitle

\noindent\emph{Source and licence.} On a Problem Posed by Brezis and Mironescu. Version arXiv:2603.07877v1 (CC BY 4.0), distributed by arXiv under the Creative Commons Attribution 4.0 International licence, http://creativecommons.org/licenses/by/4.0/, as recorded in the arXiv licence metadata for this exact version and shown on the abstract page https://arxiv.org/abs/2603.07877v1. Full licence notice: \texttt{licenses/arxiv-2603.07877-CC-BY-4.0.txt}.\par\smallskip

\noindent\textit{Editorial scope.} Counted unit: the article's proposition proposition (source0.tex lines 535--543). The article's complete original proof of it is delivered with it (source0.tex lines 544--611, one \texttt{proof} environment of 68 lines). No mathematics is written, repaired or re-derived: the statement and the proof are the article's own lines, in the article's own notation, and no retained line is substituted or re-flowed.  No further source-local context is needed: every cross-reference of every retained line resolves inside the two delivered spans.  Nothing was omitted from any retained span, so the package carries no omission record.  No retained line contains a citation, so the wrapper carries no bibliography and no bibliography entry.  No retained line contains a figure environment, an image-inclusion command or drawing code, so no image is delivered and the package declares no asset dependency.  Editorial formatting only, in addition: an AMS \texttt{amsart} preamble that restates the article's own declarations and macros used by the retained lines, namely the environments \texttt{proposition} on separate counters, together with the standard packages those lines need.  The retained environments print in this document as Proposition 1 in order of appearance, whereas the article prints the same items with the numbering of its own full document.  The complete native source of the article as distributed in the arXiv e-print for this exact version is bundled unchanged as \texttt{sources/195978/source0.tex}.
\begin{proposition} {\it Let $\Omega \subset \mathbb{R}^N$, $N\geq 3$, be a bounded convex open set. Given an integer $l$ with $0 < l < N-2$, suppose $T$ is an $l$-rectifiable current such that $T = \partial M_0$ for some smooth $(l+1)$-dimensional manifold $M_0 \subset \Omega\times\R$. Set 

\[ A^l(T) = \inf \Big\{ \mathbf{M}(\Gamma)\ | \Gamma \subset \Omega \text{ is an } (l+1)\text{-rectifiable current such that } \partial \Gamma = T \Big\},\]

\begin{align*} A_0^l(T) &= \inf \Big\{ |M|\ |\ M\subset\Omega\times\R \text{ is an } (l+1)\text{-smoothly immersed and oriented}\\
&\qquad\qquad\qquad\qquad\qquad\qquad\qquad\qquad\qquad\quad \text{submanifold with } \partial M = T \Big\}.
\end{align*}
Then $A_0^l(T) = A^l(T).$}
\end{proposition}

\begin{proof} Let $\Gamma$ be the minimal current that achieves the infimum $A^l(T)$. Let $\mathcal{S}$ be the set of singular points of $\Gamma$.
Fix $\varepsilon'>0$. Choose an $\varepsilon'$ tubular neighborhood $E_{\varepsilon'}$ of $\mathcal{S}$ given by Lemma 3.3 so that

\[\mathcal{H}^{l+1}(E_{\varepsilon'})<C \varepsilon^2,\]  
\[\mathcal{H}^{l}(\partial E')<C \varepsilon,\]
and the corresponding cone $D_\varepsilon$ given by Lemma 5.1 such that \[\mathcal{H}^{l+1}(D_\varepsilon)<C\varepsilon.\]

{To ensure that $M_0\cup \Gamma$ is immersed at the common boundary $T$ between $M_0$ and $\Gamma$, we slightly deform $M_0$ in the positive $x_{N+1}$-direction.
More precisely, let ${\rm{dist}}(\cdot, T):\R^{N+1}\to\R_+$ denotes the distance function to $T=\partial M_0$. Then there exists a $\delta_0>0$ such that 
${\rm{dist}}(\cdot, T)$ is smooth in $2\delta_0$-neighborhood
of $T$. Let $\eta\in C^\infty([0,\infty))$ satisfy $0\le\eta\le 1$, $\eta(t)=t$ for $0\le t\le \delta_0$, and $\eta(t)=2\delta_0$ for $t\ge 2\delta_0$.  Now we define
$\overline{M}_0=\big\{(x, x_{N+1}+\eta({\rm{dist}}((x, x_{N+1}),T)): (x, x_{N+1})\in M_0\big\}$. It is not hard to see that 
$$\mathcal{H}^{l+1}(\overline{M}_0)\le C\mathcal{H}^{l+1}(M_0).$$}

Furthermore, using Lemma 4.3, perform a spherical inversion of $A = \overline{M}_0 \cup T\cup \Gamma \setminus (E'_\varepsilon \cup \mathcal{S})$ through 
{the $N$-dimensional sphere $\mathbb{S}^N_r(0)$ with center $0$ and radius $r$}, choosing {$0<r<< \varepsilon^{\frac{1}{l+1}}$ 
sufficiently small  so that }the spherical inversion image $\widetilde{A}$ of $A$
has 
\[\mathcal{H}^{l+1}(\widetilde{A})\le C r^{2l+1} \mathcal{H}^{l+1}(A)\le C \varepsilon.\]

\smallskip
{To ensure the truncated cone $\mathcal{D}_\varepsilon \cup (\Gamma\setminus(E'\cup\mathcal{S}))$ is immersed at $\partial(\Gamma\setminus(E'\cup\mathcal{S}))$,
we also deform the truncated cone $\mathcal{D}_\varepsilon$ into $\widehat{\mathcal{D}_\varepsilon}$ 
by moving the vertex from $0\in\R^{N+1}$ to $(0',\varepsilon)$. More precisely, define 
\begin{align*}
\widehat{\mathcal{D}_\varepsilon}=\Big\{(tx, (1-t)\varepsilon)\ \big|\ x\in \partial E', \ \frac{r^2}{|x|^2} \le t\le 1\Big\}.
\end{align*}
It is clear that $\partial \widehat{\mathcal{D}_\varepsilon}=\partial E'\cup \big\{\big(\frac{r^2}{|x|^2} x, (1-\frac{r^2}{|x|^2})\varepsilon\big): x\in\partial E'\big\}$.
Also, it is not hard to see that there exists a smooth isometric map $\Phi:\mathbb R^{N+1}\to \mathbb R^{N+1}$ such that 
$$\Phi(\frac{r^2}{|x|^2} x)= \Big(\frac{r^2}{|x|^2} x, \ (1-\frac{r^2}{|x|^2})\varepsilon\Big), \ \forall x\in \partial E'.$$

Define $\widehat{A} =\Phi\big(\widetilde{A}\big)$. Then we have 
\begin{align*}\mathcal{H}^{l+1}\big(\Gamma \setminus (E_{\varepsilon'} \cup \mathcal{S}) \cup \widehat{D_\varepsilon} \cup \widehat{A}\big) 
&= \mathcal{H}^{l+1}(\Gamma) - \mathcal{H}^{l+1}(E_{\varepsilon'}) + \mathcal{H}^{l+1}(\widehat{D_\varepsilon}) 
+ \mathcal{H}^{l+1}(\widehat{A})\\
&\le \mathcal{H}^{l+1}(\Gamma) + C\mathcal{H}^{l+1}({D_\varepsilon}) 
+ C\mathcal{H}^{l+1}(\widetilde{A})\\
& \leq \mathcal{H}^{l+1}(\Gamma) + C \varepsilon
\end{align*}}
for some $C>0$ independent of $\varepsilon$.

Note that $\Gamma \setminus (E' \cup \mathcal{S}) \cup {\widehat{D_{\varepsilon}}
\cup \widehat{A}}$ is a finite union of smoothly immersed and oriented submanifolds.
In fact, $\Gamma \setminus (E'\cup \mathcal{S})$ is a smoothly immersed submanifold with boundary $\partial E'$, and
similarly it is true for $ \widetilde{A}$ (the image under the inversion with respect to the sphere described above) and the smoothly immersed cylinderical
surface {$\widehat{D_\varepsilon}$}. Let us denote $X$ as  an abstract manifold {of dimension $(l+1)$} and let $F : X\to \R^{N+1}$ be a smooth immersion
{whose image is} $\Gamma \setminus (E'\cup \mathcal{S})$,
namely, $dF$ has full rank {$(l+1)$} at every point in $X$. The immersion for $ \widetilde{A}$ would simply be the composition of $F$ 
{with the spherical} inversion map {described above}. With this,
the cylindrical surface {$\widehat{D_\varepsilon}$} can be viewed as a smooth immersion $G(x,t): \partial X \times [0,1]\to {\R^{N+1}}$.
Now  the rank of the differential of the map $G$  equals to  the rank of the differential of  the restriction $F\big|_{\partial X}$ plus one. Hence  $dG$ is  of full rank $(l+1)$.
%%Near each boundary point of  $\partial E'_\varepsilon$ (also a boundary point of $\partial D$) one then do a proper smoothing to link the above immersions $F$ and that for $D$.
Since the immersion $F$ {for $\Gamma\setminus (E'\cup\mathcal{S})$ and the immersion for $\widehat{D_\varepsilon}$} has the same restriction on $\partial X$, it is not hard to see one can smooth locally $F$ and that immersion for {$D_\varepsilon$ to get} a global smooth immersion of $X\cup(\partial X \times [0,1])$ into
 {$\R^{N+1}$}.  To see the latter fact, one considers $X_{\delta}$ which is the $\delta$ neighborhood of $\partial X$ in $X$. For $\delta>0$ sufficiently small,  one may view $X_{\delta}$ 
as $\partial X \times [-\delta, 0]$, and  $F$ is the smooth immersion from $X_{\delta}$  into a small open neighborhood of
$\partial E'$ into $\Gamma \setminus (E'\cup \mathcal{S}).$  For $t \in [-\delta, 0]$, 
$F_{x}$  and $F_t$ together forms a $(l+1)$ tangent frame of the image $F(X_{\delta})$, where $x \in  \partial X$
and  $l = {\rm{dim}} \ \partial E{_\varepsilon'}$.  Now what one needs to do is to extend the immersions $F$ defined on  $\partial X \times [-\delta, 0]$ and $G$ defined on $\partial X \times [\delta, 1]$ smoothly into  $\partial X \times [0, \delta]$ to become an immersion on   $\partial X \times [-\delta, 1]$ .  It is obvious what to do when $\partial X $ is a point, one simply replaces the image as a union of two smooth short arcs which intersects at a point (it could be either transversal or tangential) by a smoothly immersed slightly longer arc (in the neighborhood of the (outside) boundary points of the whole union of two arcs, one may even assume this
immersed arc are the same as the original two arcs).  For every point $x^* \in  \partial X$, one can 
{repeat what we did above}  in small neighborhood of $x^*$, and locally the resulting smooth arc changing smoothly  
when $x$ varies near $x^*$. In other words, one simply constructs a smoothly varying family (indexed by $x$) of smooth immersed arcs described as above with given given locations and its tangents at end points.  It is then by a partition of unity, one can construct such an extended  immersion on  
$\partial X \times [0, \delta]$. Similarly, one can do  smoothing near the common boundaries of {$\widehat{D_\varepsilon}$ and $\widehat{A}$}. If one does above construction with small $\delta$ so that the $(l+1)$-dimension Hausdorff measure of the image of $F (\partial X \times [0, \delta])$ is small 
(and goes to zero as $\delta$ goes to zero, and in particular much less than $\varepsilon$), then
we call the resulting approximation by $M_\varepsilon$ so that
\[ A^l_0(T) \leq \mathcal{H}^{l+1}(M_\varepsilon) \leq \mathbf{M}(\Gamma) + C\varepsilon,\]
noting that $M_\varepsilon$ satisfies $\partial M_\varepsilon= T$.
Then we can not have $A^l(T) < A_0^l(T)$, but we {have already shown $A^l(T) \leq A_0^l(T)$. Therefore,} $A^l_0(T)=A^l(T).$
\end{proof}

\end{document}
